\scanpage{104}
\uline{Définition} 1\quad Soit $X$ un espace topologique. On appelle topologie constructible de $X$, \deleted{\unsure{dont}} la topologie dont les ouverts sont les réunions quelconques de parties loc. constructibles [donc une base d'ouverts fournie par les parties loc. constructibles,]. Parties \unsure{antifiques|antifinies} $=$ fermés de $X^{\mathrm{cons}}$ $=$ intersections de parties loc. constructibles.

\margin{\deleted{\unsure{Universel} \illegible{2}}}
\uline{Prop 1}\quad Supposons que \deleted{tout point de} $X$ admette une base \inserted{$(U_i)_{i\in I}$} formée d'ouverts rétrocompacts. Alors

(0) les parties constructibles de $X$ forment une base d'ouverts de $X^{\mathrm{cons}}$

i) La topologie constructible de $X$ est plus fine que la topologie donnée, i.e. $X^{\mathrm{cons}}\unsure{\text{---}|\longrightarrow}X$ est continue.

ii) \deleted{\unsure{Si on}}\inserted{Si} $X$ satisfait de plus $T_0$, alors $X^{\mathrm{cons}}$ est séparé [$=$ tot. discontinu] et réciproquement.

iii) Si $\unsure{E}$ est une \deleted{ouvert de $X$, la topologie}\inserted{partie constructible de $X$,} munie de sa topologie induite $T_E$, alors la topologie constructible $T_E^{\mathrm{cons}}$ associée à $T_E$ est induite par $T^{\mathrm{cons}}$ sur $E$.

\margin{\figure{\begin{tikzpicture}[x=1cm,y=1cm]
\draw (0,0)--(0,-0.65);
\draw (0.20,-1.0)--(0.24,-1.6);
\draw (0.34,-1.95)--(0.42,-2.45);
\draw (0.48,-2.85)--(0.50,-4.15);
\end{tikzpicture}}}

(0) Car un loc. constructible $E$ est réunion des $U_i\cap E$, qui est constructible dans $X$ par $\unsure{0_{\mathrm{III}}|0_{\mathrm{II}}}$ 9.1.8.

i) Résulte du fait que les $U_i$ sont \inserted{[loc]} constructibles, donc ouverts pour $T_{\unsure{X}}^{\mathrm{cons}}$. Pour ii), \deleted{\unsure{soit}}, $x,y\in X$, il existe un \unsure{des deux, soit} $x$, et un voisinage ouvert $U\ni x$, $U\not\ni y$. On peut grâce à l'hyp supposer $U$ rétrocompact, donc [loc.] constructible, et $X-U=V$ l'est \inserted{alors} également, alors
\[
U\cap V=\varnothing,\quad U\cup V=X,\quad U\ni x,\quad V\ni y,\quad U,V\text{ ouverts pour }T_{\unsure{E|X}}.
\]

Pour iii), on sait que la topologie induite par $T_{\unsure{X|E}}^{\mathrm{cons}}$ sur $E$ a comme base d'ouverts les $E\cap F$, $F$ \deleted{\unsure{co}}\inserted{loc}. constructible dans $X$, ce sont exactement les $G\subset E$ loc. constructibles dans $E$,

\margin{103}
